\section{Warm start 证据：曲线应该怎样读}

课程用两组实验说明初始化可能跨任务有用：单层语言模型的困惑度轨迹，以及 MNIST 的单层/双层分类。图表的重点是起点和优化路径，而不是宣布新架构已经达到大模型质量。

\subsection{单层语言模型：起点更低，轨迹仍继续改善}

前一章只给出了参数计算方法，本节转向最直接的经验问题：这个 data-informed 初值是否真的比随机起点更接近可用模型？读曲线时先比较 epoch 0 的起点，再比较相同 update rule 下的下降轨迹与 early-stopping 位置；不要先看最终一个点，就忽略 warm-start statistics 本身也有计算成本。

Perplexity（困惑度）是平均 cross-entropy 的指数：
\[
\operatorname{PPL}=\exp\!\left(-\frac{1}{T}\sum_{t=1}^{T}\log p(x_t\mid x_{<t})\right).
\]
它可直观理解为模型每个 token 面对的“有效候选数”，越低通常越好；但不同 tokenization、数据和 context 设置的 PPL 不能直接横比。

\lecturefigure{slide-09-single-layer-warm-starts.jpg}{单层模型中 warm start 与 cold start 的训练轨迹}{Stanford Online 官方录像 00:24:03}

读图先看 epoch 0：cold start 的 PPL 接近 vocabulary size，warm start 已明显降低。再看斜率和终点：在相同 iterative update 规则下，warm-start 曲线继续下降，并在早期停止时保持更低 PPL。它支持“显式初值改善起点和后续优化”，但没有控制所有现代 Transformer 组件。

\begin{knowledgebox}{读图：这张图没有 self-attention}
Q\&A 明确说明，这是 concatenated one-hot context 的 feed-forward experiment。输入 feature 已分开，几乎不需要 attention 去解除 superposition。它验证的是 single-layer explicit initialization，不是 SAFFU attention 的完整消融。
\end{knowledgebox}

\begin{warningbox}{同一 learning rate 不等于完全公平}
Warm start 和 random start 可能适合不同 learning-rate schedule、regularization 或 early-stopping patience。严谨比较应同时报告 tuned cold baseline、多个 seeds、wall-clock、数据读取与显式统计量计算成本。
\end{warningbox}

\teachervoice{00:22:47--00:24:24，讲者强调 warm start 在训练前就降低 PPL，之后还沿相同 learning rate 继续改善；若使用 early stopping，也会更早在更低 PPL 处触发。}

\teachervoice{01:07:53--01:09:37，Q\&A 纠正了常见误读：这张 warm/cold 图没有 attention，早期工作只是单层 feed-forward，SAFFU attention 的显式解是后来才补上的。}

\subsection{MNIST：说明 layer-level 公式不只适用于文字}

MNIST pixel 非负，因此同一 softmax layer warm start 可以作用于 image classification。图中 warm-start 单层和双层模型更快达到高 accuracy；但模型没有现代视觉网络的 convolution、normalization 或 augmentation pipeline。

\lecturefigure{slide-10-image-classification-warm-starts.jpg}{MNIST 上的 warm-start 分类曲线}{Stanford Online 官方录像 00:30:45}

\begin{importantbox}{MNIST 结果支持的最小结论}
显式 softmax initialization 可以接受非语言、非负 feature，并改善该分类设置中的起点。它不证明同一公式已经覆盖 convolution、ViT、diffusion 或真实世界视觉数据。
\end{importantbox}

\teachervoice{00:29:14--00:31:20，讲者用 MNIST 说明公式的 layer-level generality；他没有声称已经构造完整高性能 vision architecture。}

\teachervoice{01:09:42--01:11:46，Q\&A 进一步限制范围：若要处理 convolution 或其他 activation，仍需为那些 layer 单独推导 warm start，否则随机参数会重新引入噪声。}

\subsection{实验应该补哪些 baseline}

对于这种 initialization 研究，单看最终 accuracy/PPL 不够。至少要拆出四类成本：显式统计量计算、随机 baseline 的 tuning、后续 backpropagation、以及固定/更新 embedding 对缓存的影响。

\begin{knowledgebox}{最小实验矩阵}
\begin{tabularx}{\textwidth}{L{0.18\textwidth} L{0.24\textwidth} X}
\toprule
对比 & 固定项 & 需要报告 \\
\midrule
Cold vs warm & architecture、data、optimizer & seeds、初始/最终 loss、wall-clock \\
Warm only vs warm+tune & initialization statistics & 额外梯度成本与泛化增益 \\
Frozen vs thawed embeddings & warm-start weights & cache 命中、训练速度、最终质量 \\
Local explicit vs end-to-end & parameter count & compositional gap 与 failure cases \\
\bottomrule
\end{tabularx}
\end{knowledgebox}

\subsection{本章小结}

Warm start 实验显示更好的起点和优化轨迹，并在 MNIST 上展现有限跨域性。它们验证的是特定 layer 和数据条件，不是“所有深度学习都可被共现矩阵替代”。

